% Chapter 01. Source: import/awyi_optimisation.tex (18 Nov 2026), ported in D.2 without rewriting:
% title block and table of contents removed (the subtitle, byline and date are the \sourcenote),
% macros moved to preamble.tex, \Lh renamed \Lha (fpo_from_zero.tex uses \Lh for \hat L), labels
% prefixed with loop:, the TikZ picture in figures/01_tensor.tex, and the appendix
% (A worked example) kept as the last section.
\chapter{Achieve What You Imagined: the optimisation, step by step}\label{ch:paper-loop}
\sourcenote{Rewards, returns, the $(b,i,k)$ ``sandbox'', and the FPO update of the action head.\\
Study notes on arXiv:2609.33832, Sections IV.D--IV.E and V.A, cross-checked against the \texttt{imagine-to-achieve} code. 18 November 2026.}
\begin{abstract}
These notes rebuild the training procedure of the paper from the reward to the weight update, filling the steps the paper compresses. Paper equation numbers are written as \eqp{8}--\eqp{16}; equations of these notes use plain numbers. Three things are worked out in detail: (i) how the two rewards become one number per chunk and then a discounted suffix return (Eq.~\eqp{11}); (ii) how the group normalisation \eqp{12} restricts every comparison to a small set of rollouts with the same start and the same chunk position, and what that does and does not cancel; (iii) how the clipped FPO objective \eqp{13}--\eqp{16} turns those advantages into gradients on the action head only. A final section collects the experimental setup, including what the paper does not report. Where the code differs from or adds to the paper, a grey box says so; those boxes are based on reading the synced source files, not on running them, and the YAML configs (where the actual values of $Q$, $\epsilon_{\mathrm{clip}}$, the learning rate and $\lambda$ live) are not in the synced folders.
\end{abstract}

%=====================================================================
\section{The objects and their indices}
%=====================================================================

\subsection{One chunk of one rollout}
At chunk $k$ the world-action model (Cosmos3) samples a joint latent endpoint with its native UniPC solver and decodes it:
\begin{equation}
z_{b,i,k}=\big[z^{v}_{b,i,k};\,z^{a}_{b,i,k}\big]\sim \Pi^{z}_{\theta}(\cdot\mid o_{b,i,k},g_b),\qquad
\big(\VC_{b,i,k},\,a_{b,i,k}\big)=D(z_{b,i,k}).
\end{equation}
$\VC$ is the imagined plan video (``what the policy thinks will happen''), $a\in\R^{H_a\times d_a}$ the action chunk ($H_a=32$ steps at 15\,Hz, $d_a=10$: body-frame position delta, rot6d rotation delta, absolute gripper). Ctrl-World, frozen, receives the converted actions and predicts what they produce:
\begin{equation}
\VW_{b,i,k}=W_{\phi^0}\big(o_{b,i,k},T_a(a_{b,i,k})\big),\qquad o_{b,i,k+1}=\mathrm{Ext}_{\mathrm{obs}}\big(\VW_{b,i,k}\big).
\end{equation}
The next observation comes from $\VW$, never from $\VC$. From chunk~1 on, the policy therefore only sees images generated by Ctrl-World.

\subsection{The four indices}
\begin{center}
\small
\begin{tabularx}{\linewidth}{@{}l l Y@{}}
\toprule
Index & Range (paper) & Meaning and what is \emph{shared} along it\\
\midrule
$b$ & $1,\dots,B$, $B=8$ & \textbf{group}. One reset: a start observation $o_{b,0}$ (first frame of a demonstration) and an instruction $g_b$. Different $b$ = different scenes.\\
$i$ & $1,\dots,G$, $G=16$ & \textbf{member} of a group. All $i$ in group $b$ share $o_{b,0}$ and $g_b$; they differ only through the policy's sampling noise (and Ctrl-World's).\\
$k$ & $0,\dots,N-1$, $N=5$ & \textbf{chunk position} along a rollout. One chunk = one policy call = 32 actions.\\
$t$ & $1,\dots,H_v$ & \textbf{frame} inside a chunk, at the timestamps shared by $\VC$ and $\VW$ (only used inside the reward).\\
$q$ & $1,\dots,Q$ & \textbf{probe} (time--noise pair) used to evaluate the FPO loss of one stored chunk.\\
\bottomrule
\end{tabularx}
\end{center}

One update therefore works on a tensor of $B\times G\times N=8\times16\times5=640$ chunk samples. \Cref{loop:fig:tensor} shows the layout and the two operations that act on it: the reward and return run \emph{along a row} (fixed $b,i$, over $k$); the normalisation runs \emph{down a column inside one block} (fixed $b,k$, over $i$). Everything below is about these two directions.

\begin{figure}[t]
\centering
% The figure "loop:fig:tensor" of import/awyi_optimisation.tex, moved here unchanged in D.2;
% \input by chapters/01_paper_loop.tex. The final % keeps the spacing of the original.
\begin{tikzpicture}[x=0.78cm,y=0.5cm,>={Stealth[length=2mm]}]
  \foreach \b/\yo/\lab in {0/0/{b=1},1/-6.2/{b=2},2/-12.4/{b=B}}{
    \foreach \i in {0,...,4}{
      \foreach \k in {0,...,4}{
        \pgfmathsetmacro{\yy}{\yo-\i}
        \draw[gray!60,fill=white] (\k,\yy) rectangle ++(1,-0.9);
      }
    }
    \node[anchor=east,font=\small,align=right] at (-0.35,\yo-2.4) {$\lab$\\[-1pt]{\scriptsize members}\\[-2pt]{\scriptsize $i=1,\dots,G$}};
    \draw[decorate,decoration={brace,amplitude=4pt,mirror}] (-0.1,\yo) -- (-0.1,\yo-4.9);
  }
  \node at (2.5,-11.4) {$\vdots$};
  \draw[very thick,blue!70!black,fill=blue!12] (2,0) rectangle (3,-4.9);
  \foreach \i in {0,...,4}{ \draw[gray!60] (2,-\i) rectangle ++(1,-0.9); }
  \draw[very thick,red!70!black] (0,-1) rectangle (5,-1.9);
  \foreach \k in {0,...,4}{ \node[font=\scriptsize] at (\k+0.5,0.6) {$k{=}\k$}; }
  \node[align=left,font=\small,anchor=north west,text width=6.4cm] (rowlab) at (6.6,1.2)
    {\textcolor{red!70!black}{\textbf{Row} $(b,i,\cdot)$: one rollout.} Rewards $r_{b,i,k}$ are computed chunk by chunk and summed backwards into the suffix return $R_{b,i,k}$, Eq.~\eqp{11}.};
  \node[align=left,font=\small,anchor=north west,text width=6.4cm] (collab) at (6.6,-4.4)
    {\textcolor{blue!70!black}{\textbf{Column inside a block} $(b,\cdot,k)$: the comparison set $\mathcal S_{b,k}$.} Mean and std are taken over these $G$ cells only, Eq.~\eqp{12}. Cells in other blocks or other columns never enter.};
  \draw[->,red!70!black] (6.55,-1.45) -- (5.05,-1.45);
  \draw[->,blue!70!black] (6.55,-5.6) -- (3.05,-3.8);
  \node[align=left,font=\small,anchor=north west,text width=6.4cm] at (6.6,-11.0)
    {After normalisation all $BGN$ cells are flattened into one list of chunk samples $(z,c,\Omega,\hat L^{a}_{\mathrm{old}},\hat A)$, split into minibatches and used by the FPO update, Eq.~\eqp{16}. From then on the indices only label samples.};
\end{tikzpicture}%

\caption{The $B\times G\times N$ batch of one update (only 5 of the 16 members are drawn per group).}
\label{loop:fig:tensor}
\end{figure}

%=====================================================================
\section{From two rewards to one number per chunk}
%=====================================================================

\subsection{Trajectory reward \eqp{8}}
For each chunk the two videos are sampled at $H_v$ common timestamps and compared view by view at the same resolution:
\begin{equation}
r^{\mathrm{traj}}_{b,i,k}=-\frac1{H_v}\sum_{t=1}^{H_v}D\big(\VC_{b,i,k,t},\VW_{b,i,k,t}\big),\qquad
D(F,F')=\frac{\|F-F'\|_2^2}{P},
\label{loop:eq:traj}
\end{equation}
with pixels on the $[0,255]$ scale and $P$ the number of compared values (channels $\times$ views $\times$ pixels). It asks: \emph{did the actions do what the policy imagined they would do?} Note that the target $\VC_{b,i,k}$ is the member's \emph{own} plan, so two members are each scored against a different target.

\subsection{Goal reward \eqp{9}}
At the last chunk only, the final Ctrl-World frame is compared with a fixed goal image (from the demonstration):
\begin{equation}
r^{\mathrm{goal}}_{b,i}=-D\big(\VW_{b,i,N-1,\mathrm{final}},\,I^\star_{g}\big).
\end{equation}
This target is \emph{the same for all members of a group}, unlike the trajectory target.

\subsection{Combined reward \eqp{10}}
\begin{equation}
r_{b,i,k}=\lambda_{\mathrm{traj}}\,r^{\mathrm{traj}}_{b,i,k}+\lambda_{\mathrm{goal}}\,\mathbf 1\{k=N-1\}\,r^{\mathrm{goal}}_{b,i}.
\label{loop:eq:reward}
\end{equation}
All rewards are $\le0$; ``higher'' means ``smaller pixel error''. The reward is a plain number: it is computed under \texttt{no\_grad}, and no gradient ever flows from it into any network. This is the first sense of the sandbox: the policy cannot change the judge (both video models are frozen), and it cannot even differentiate through the judge; it can only change which actions it samples.

\begin{codenote}[(rewards)]
\begin{itemize}[leftmargin=*]
\item The training reward is selected by \texttt{reward.training\_source}: \texttt{continuous\_combined} (\eqref{loop:eq:reward}), \texttt{trajectory\_mse}, \texttt{terminal\_goal\_mse} (the two ablation arms of Fig.~6), or \texttt{success\_binary} (a classifier reward). In the paper runs the success classifier is \emph{diagnostic only}.
\item The plotting manifest of a Push-T run lists per-camera weights for the trajectory term: wrist $2/3$, main $1/6$, side $1/6$ (and a ``no\_wrist'' variant with $0,\tfrac12,\tfrac12$). So in practice the views are \emph{not} pooled with equal weight as ``$P$ counts all values'' suggests. Check the config of the task you reproduce.
\end{itemize}
\end{codenote}

%=====================================================================
\section{Credit along a rollout: the suffix return \eqp{11}}
%=====================================================================

\subsection{Definition and matrix form}
\begin{equation}
R_{b,i,k}=\sum_{j=k}^{N-1}\gamma^{\,j-k}\,r_{b,i,j},\qquad \gamma=0.95.
\label{loop:eq:return}
\end{equation}
For $N=5$, writing the row $(b,i)$ as vectors,
\begin{equation}
\begin{pmatrix}R_0\\R_1\\R_2\\R_3\\R_4\end{pmatrix}
=\underbrace{\begin{pmatrix}
1&\gamma&\gamma^2&\gamma^3&\gamma^4\\
0&1&\gamma&\gamma^2&\gamma^3\\
0&0&1&\gamma&\gamma^2\\
0&0&0&1&\gamma\\
0&0&0&0&1
\end{pmatrix}}_{\Gamma}
\begin{pmatrix}r_0\\r_1\\r_2\\r_3\\r_4\end{pmatrix},
\qquad R=\Gamma r .
\end{equation}
The code computes this backwards, $R_{N-1}=r_{N-1}$, $R_k=r_k+\gamma R_{k+1}$ (a \texttt{done} flag would cut the recursion; with fixed-length rollouts it never does).

\subsection{What each chunk is credited for}
Inserting \eqref{loop:eq:reward} splits the return into the two sources:
\begin{equation}
R_{b,i,k}=\lambda_{\mathrm{traj}}\underbrace{\sum_{j=k}^{N-1}\gamma^{j-k}r^{\mathrm{traj}}_{b,i,j}}_{\text{own and later consistency}}
\;+\;\lambda_{\mathrm{goal}}\,\underbrace{\gamma^{\,N-1-k}\,r^{\mathrm{goal}}_{b,i}}_{\text{final outcome}} .
\label{loop:eq:split}
\end{equation}
\begin{itemize}
\item \textbf{Only the future counts.} Chunk $k$ is not credited for anything that happened before it (``suffix''), which is right: its actions cannot change the past.
\item \textbf{Every chunk sees the goal.} The goal reward reaches chunk $k$ with weight $\gamma^{N-1-k}$: $0.815,\,0.857,\,0.903,\,0.95,\,1$ for $k=0,\dots,4$. With $\gamma=0.95$ and $N=5$ the discount barely matters; it is closer to ``sum the rest of the rollout'' than to a strong preference for immediate reward.
\item \textbf{Different chunk positions have different scales.} $R_0$ sums five terms, $R_4$ only one. In the paper's laptop ablation the combined per-rollout reward is around $-12$ (Fig.~6), so $R_0$ and $R_4$ live on clearly different ranges. This is one of the reasons for normalising per chunk position (\cref{loop:sec:sandbox}).
\item \textbf{Credit is not causal.} A chunk is credited with later rewards that may be mostly due to later chunks' actions (or Ctrl-World noise). There is no critic to separate ``my contribution'' from ``what came after''; the group comparison is the only baseline.
\end{itemize}

%=====================================================================
\section{The sandbox: group-relative normalisation \eqp{12}}
\label{loop:sec:sandbox}
%=====================================================================

\subsection{Definition}
For each group $b$ and chunk position $k$, define the comparison set $\mathcal S_{b,k}=\{(b,i,k):i=1,\dots,G\}$ and
\begin{equation}
\Ah_{b,i,k}=\frac{R_{b,i,k}-\mu_{b,k}}{s_{b,k}+\delta},\qquad
\mu_{b,k}=\frac1G\sum_{i=1}^{G}R_{b,i,k},\quad
s_{b,k}=\Big(\frac1G\sum_{i=1}^{G}(R_{b,i,k}-\mu_{b,k})^2\Big)^{1/2}.
\label{loop:eq:adv}
\end{equation}
There are $B\cdot N=40$ such sets per update, each of size $G=16$. Every quantity that enters the gradient is computed \emph{inside} one set; nothing is compared across sets. This replaces the learned value function (critic) of PPO: the baseline for a chunk is ``how the other 15 members did at the same position from the same start''.

\begin{codenote}[(advantages)]
The helper \texttt{grpo\_action\_suffix\_advantages} (\texttt{src/rlinf\_modified/algorithms/advantages.py}) uses \texttt{torch.std}, which by default is the \emph{Bessel-corrected} std (divide by $G-1$), with $\delta=10^{-6}$; the implementation notes for the production path (\texttt{cross\_rank\_grpo.py}) also find that the sum of squared advantages is $G-1$. The paper says ``population'' std (divide by $G$). The difference is a constant factor $\sqrt{G/(G-1)}\approx1.03$ on all advantages, harmless but worth knowing when you reproduce numbers. Groups are spread over several GPUs (ranks), so rewards are all-gathered before \eqref{loop:eq:adv} is evaluated.
\end{codenote}

\subsection{What the sandbox cancels}
\begin{proposition}[Offsets cancel within a set]\label{loop:prop:shift}
Let $c_{b,k}$ be any numbers. Replacing $r_{b,i,k}$ by $r_{b,i,k}+c_{b,k}$ for every member $i$ leaves every advantage $\Ah_{b,i,k'}$ unchanged, for all $k'$.
\end{proposition}
\begin{proof}
By \eqref{loop:eq:return} the shift adds $\gamma^{k-k'}c_{b,k}$ to $R_{b,i,k'}$ for $k'\le k$ (and nothing for $k'>k$). The added amount does not depend on $i$, so $\mu_{b,k'}$ moves by the same amount and both $R_{b,i,k'}-\mu_{b,k'}$ and $s_{b,k'}$ are unchanged.
\end{proof}
\textbf{Consequences.}
\begin{itemize}
\item \emph{Scene difficulty cancels.} If one start scene (group $b$) gives large pixel errors for everyone (clutter, lighting, a particular duck colour), all its members are shifted together and none of them is pushed. Without grouping, every member of an easy scene would get a positive advantage and the policy would learn ``do what I did on easy scenes''.
\item \emph{Chunk-position effects cancel.} Effects common to a position (later chunks drift more because Ctrl-World conditions on its own predictions; only the last chunk carries the goal term) are removed because $k$ is never mixed.
\item \emph{The goal image's absolute distance does not matter.} Only differences between members' final frames do.
\end{itemize}

\begin{proposition}[Fixed budget per set]\label{loop:prop:budget}
If $s_{b,k}\gg\delta$, then $\sum_{i}\Ah_{b,i,k}=0$ and $\sum_i\Ah_{b,i,k}^2=G$ (or $G-1$ with Bessel's correction), independently of the rewards.
\end{proposition}
\begin{proof}
Centring gives zero sum; dividing by the std gives unit mean square.
\end{proof}
This is the ``dilution'' you noticed, and it has two faces:
\begin{itemize}
\item \textbf{Good:} every set pushes with the same strength, so a task or a scene with large reward magnitudes does not dominate the update. Multiplying all rewards by a constant $c>0$ changes nothing (up to $\delta$), so only the \emph{ratio} $\lambda_{\mathrm{goal}}/\lambda_{\mathrm{traj}}$ matters, not their overall size.
\item \textbf{Bad:} a set whose returns differ only by noise (e.g.\ Ctrl-World sampling noise, or $16$ members that all fail in nearly the same way) is pushed exactly as hard as a set with real differences. Normalisation amplifies small differences to unit scale. $\delta$ only protects against $s=0$, not against $s$ being small and meaningless. This is the reason why your thesis's referee/weighting has to think about scale (decision D1 in the implementation plan).
\end{itemize}

\begin{proposition}[Which reward term ``wins'']\label{loop:prop:mix}
Write $R_{b,i,k}=\lambda_{\mathrm{traj}}T_{b,i,k}+\lambda_{\mathrm{goal}}\gamma^{N-1-k}Y_{b,i}$ (with $Y_{b,i}=r^{\mathrm{goal}}_{b,i}$) as in \eqref{loop:eq:split}. Then
\[
\Ah_{b,i,k}=\frac{\lambda_{\mathrm{traj}}\,(T_{b,i,k}-\bar T_{b,k})+\lambda_{\mathrm{goal}}\gamma^{N-1-k}\,(Y_{b,i}-\bar Y_{b})}{s_{b,k}+\delta},
\]
so the ranking within a set is driven by $\lambda_{\mathrm{traj}}\,\mathrm{std}_i(T)$ versus $\lambda_{\mathrm{goal}}\gamma^{N-1-k}\,\mathrm{std}_i(Y)$, i.e.\ by the \emph{spread} of each term inside the group, not by its average size.
\end{proposition}
So setting $\lambda_{\mathrm{traj}}=\lambda_{\mathrm{goal}}$ does not mean ``equal importance''. If the members' final frames differ a lot (some closed the lid, some did not) while their trajectory consistency is similar, the goal term decides the advantage, and vice versa. The worked example in \cref{loop:app:toy} is exactly such a case.

\subsection{What the sandbox does \emph{not} control}
\begin{caveat}[Only chunk 0 is a clean same-state comparison]
All members share $o_{b,0}$, but $o_{b,i,1}=\mathrm{Ext}_{\mathrm{obs}}(\VW_{b,i,0})$ already differs between members. So for $k\ge1$ the set $\mathcal S_{b,k}$ compares chunks taken \emph{from different states}, with a baseline $\mu_{b,k}$ that is the same for all of them. A member that is in a worse state at $k=2$ (because of chunk 0 or 1) gets a negative advantage at $k=2$ even if its chunk-2 action was the best available from where it was. The return-to-go already carries chunk 0's consequences, so this mainly adds noise to the later-chunk updates rather than bias in the sign; still, the cleanest signal sits at $k=0$. This is the reason the thesis proposal restricts its interventional checks to chunk~0.
\end{caveat}
\begin{itemize}
\item The sandbox removes offsets \emph{common to a group}. It does not remove errors that differ between members, for instance Ctrl-World being right about one member's contact and wrong about another's. Those flip rankings, and rankings are all the gradient sees.
\item The trajectory term uses each member's \emph{own} plan as target. A member with a modest plan that is executed faithfully can outrank a member with an ambitious plan that is executed partly. The goal term is what keeps the policy from ``planning small''.
\end{itemize}

%=====================================================================
\section{Section IV.E in detail: FPO on the action head}
%=====================================================================

\subsection{Why not a normal policy gradient}
A policy gradient needs $\nabla_\theta\log\pi_\theta(a\mid c)$. For a flow model the likelihood of a sample needs the full ODE (and a trace of the Jacobian), and for UniPC, a multi-step high-order solver, there is not even a simple per-step Gaussian to write down (which Flow-GRPO and $\pi$RL need). FPO avoids all of this: it scores the \emph{endpoint} $z$ with a conditional flow-matching (CFM) loss and treats the sampler as a black box. Only the stored endpoint matters, not how UniPC reached it.

\subsection{Noising the stored endpoint \eqp{13}}
For each stored chunk, draw $Q$ probes $\Omega_{b,i,k}=\{(\tau_q,\varepsilon_q)\}_{q=1}^Q$, $\tau_q\sim\mathcal U(0,1)$, $\varepsilon_q\sim\mathcal N(0,I)$, and form
\begin{equation}
z^{(q)}=z+(1-\tau_q)\,M\odot(\varepsilon_q-z)\qquad(\text{indices }b,i,k\text{ dropped}).
\label{loop:eq:noise}
\end{equation}
At $\tau=1$ this is the clean endpoint $z$; at $\tau=0$ the masked coordinates are pure noise. Differentiating in $\tau$ gives the target velocity
\begin{equation}
\frac{\partial z^{(q)}}{\partial\tau}=M\odot(z-\varepsilon_q)=:u^{(q)},
\label{loop:eq:target}
\end{equation}
the constant ``noise-to-data'' direction of the straight path. The mask $M$ selects \textbf{both} the future video latents and the valid action coordinates; the conditioning latents (current frame) and padding stay clean. The video part is noised too because during sampling Cosmos3 denoises video and action \emph{jointly} at the same noise level; asking the action head for a velocity while the video latents were clean would query it in a situation it never sees.

\begin{codenote}[(time convention, noise, schedule)]
\begin{itemize}[leftmargin=*]
\item The code uses $\sigma=1-\tau$: $x_\sigma=x_0+\sigma(\varepsilon-x_0)$, target $\varepsilon-x_0$ ($\sigma=1$ is noise). Same path, reversed clock; signs of the target flip accordingly.
\item $\tau$ is not necessarily uniform: \texttt{fpo.time\_distribution} and a \texttt{training\_shift} $s$ map a base time $t$ to $\sigma=st/(1+(s-1)t)$ (the usual ``shifted'' schedule of video diffusion), which concentrates probes at high noise when $s>1$.
\item Probes are not stored as tensors but as integer seeds derived from the sampling seed and the probe index; $\varepsilon_q$ is regenerated from them. Before the first optimiser step the code recomputes $\Lha_{\mathrm{old}}$ and raises an error if it does not match the stored value. This guarantees the ``same probes for old and new policy'' condition.
\end{itemize}
\end{codenote}

\subsection{The action-only CFM loss \eqp{14}}
\begin{equation}
\Lha_{\tha}(z\mid c)=\frac1Q\sum_{q=1}^{Q}\Big\|P_a\big[v_{\thv,\tha}(z^{(q)},\tau_q\mid c)-u^{(q)}\big]\Big\|_2^2,\qquad c=(o_{b,i,k},g_b).
\label{loop:eq:cfm}
\end{equation}
$P_a$ keeps only the valid action coordinates of the predicted and target velocity. Read it as: \emph{how well does the current network ``explain'' this particular action as a denoising target, in this state?} A small $\Lha$ means that from noised versions of $z^a$ the network points straight back to $z^a$, so sampling in the same condition is likely to land near $z^a$.

Two things about the gradient:
\begin{itemize}
\item The velocity at the action coordinates depends on both parameter blocks (the action head attends to the video latents through the shared network), but only $\tha$ receives gradients; $\thv$ is frozen (\texttt{fpo.freeze\_video\_path}, \texttt{fpo.train\_action\_path\_only}).
\item The loss on video coordinates is not part of the objective (the code has a diagnostic \texttt{fpo\_vision\_velocity\_mse} that is explicitly excluded from the loss).
\end{itemize}

\begin{codenote}[(score parameterisation)]
\texttt{fpo.score\_parameterization} $\in\{\texttt{velocity},\texttt{epsilon}\}$. \eqref{loop:eq:cfm} is the velocity form. The epsilon form is the error of the implied noise prediction, which for the linear path is exactly $(1-\sigma)^2=\tau^2$ times the velocity error. It therefore reweights the probes: noisy probes ($\tau$ small) count much less. This changes how fast the ratio below moves and which probes dominate it, so it is a real design choice, not a detail. The code uses ``score'' $=-\Lha$, and \texttt{log\_ratio = current\_score $-$ old\_score}.
\end{codenote}

\subsection{The surrogate ratio \eqp{15} and why it is a ratio}
\begin{equation}
\hat\rho_{b,i,k}(\tha)=\exp\!\big(\Lha_{\thao}(z_{b,i,k}\mid c_{b,i,k})-\Lha_{\tha}(z_{b,i,k}\mid c_{b,i,k})\big),
\label{loop:eq:ratio}
\end{equation}
with the same endpoint, condition and probes in both losses. The justification (from the FPO paper) is that, for a suitable weighting over $\tau$, the CFM loss is, up to an additive constant and a positive factor, the negative evidence lower bound (ELBO) of the sample: $\log p_\theta(z\mid c)\gtrsim-\kappa\,L_{\mathrm{CFM},\theta}(z\mid c)+\text{const}$. The difference of two losses at the same $z$ is then a proxy for $\log p_\theta-\log p_{\mathrm{old}}$, and the constant cancels. Properties worth keeping in mind:
\begin{itemize}
\item $\hat\rho>1$ $\Leftrightarrow$ the loss on this sample went \emph{down} $\Leftrightarrow$ the new policy ``likes'' this action more.
\item It is a surrogate, not the true likelihood ratio: the bound is not tight, the weighting in practice is chosen empirically (\texttt{score\_parameterization}, schedule shift), and with finite $Q$ it is a Monte Carlo estimate. Sharing the probes between old and new (common random numbers) is what makes the \emph{difference} low-variance even when each loss is noisy.
\item The scale of $\Lha$ sets how fast $\hat\rho$ moves per parameter change: a sum over the $32\times10=320$ action coordinates moves it much faster than a mean. This interacts with the clip range like a learning rate does.
\end{itemize}

\subsection{The clipped objective \eqp{16}}
\begin{equation}
J_{\mathrm{FPO}}(\tha)=\frac1{BGN}\sum_{b,i,k}\min\Big(\hat\rho_{b,i,k}\Ah_{b,i,k},\ \clip(\hat\rho_{b,i,k},1-\epsilon_{\mathrm{clip}},1+\epsilon_{\mathrm{clip}})\,\Ah_{b,i,k}\Big).
\label{loop:eq:fpo}
\end{equation}
This is the PPO clip with the ratio replaced. Case by case:
\begin{center}
\small
\begin{tabular}{@{}llll@{}}
\toprule
Advantage & Ratio & Term used & Gradient\\
\midrule
$\Ah>0$ (better than group) & $\hat\rho\le1+\epsilon$ & $\hat\rho\Ah$ & pushes $\Lha$ down (make action more likely)\\
$\Ah>0$ & $\hat\rho>1+\epsilon$ & $(1+\epsilon)\Ah$ & \textbf{zero}: already moved enough\\
$\Ah<0$ (worse than group) & $\hat\rho\ge1-\epsilon$ & $\hat\rho\Ah$ & pushes $\Lha$ up (make action less likely)\\
$\Ah<0$ & $\hat\rho<1-\epsilon$ & $(1-\epsilon)\Ah$ & \textbf{zero}: already moved enough\\
\bottomrule
\end{tabular}
\end{center}

\subsection{The gradient, written out}
Since $\nabla\hat\rho=\hat\rho\,\nabla\log\hat\rho=-\hat\rho\,\nabla\Lha_{\tha}$,
\begin{equation}
\nabla_{\tha}J_{\mathrm{FPO}}=-\frac1{BGN}\sum_{b,i,k}\mathbf 1\{\text{not clipped}\}\;\Ah_{b,i,k}\,\hat\rho_{b,i,k}\;\nabla_{\tha}\Lha_{\tha}(z_{b,i,k}\mid c_{b,i,k}).
\label{loop:eq:grad}
\end{equation}
At the first minibatch of an update $\tha=\thao$, so $\hat\rho=1$, nothing is clipped, and
\begin{equation}
\nabla_{\tha}J_{\mathrm{FPO}}\big|_{\tha=\thao}=-\frac1{BGN}\sum_{b,i,k}\Ah_{b,i,k}\,\nabla_{\tha}\Lha_{\tha}(z_{b,i,k}\mid c_{b,i,k}).
\label{loop:eq:grad1}
\end{equation}
\begin{intuition}{The update is signed, advantage-weighted flow matching on the policy's own samples}
Compare with supervised fine-tuning, where the step is $-\eta\,\frac1n\sum\nabla L_{\mathrm{CFM}}(\text{demo action})$. Equation~\eqref{loop:eq:grad1} is the same regression, but (i) on actions the policy generated itself, in states Ctrl-World generated, and (ii) each sample is weighted by its advantage: positive weight $=$ ``regress towards this action'', negative weight $=$ ``regress away from it''. Because $\sum_i\Ah_{b,i,k}=0$ in every set (\cref{loop:prop:budget}),
\[
\sum_{i}\Ah_{b,i,k}\nabla\Lha_{b,i,k}=\sum_i\Ah_{b,i,k}\big(\nabla\Lha_{b,i,k}-\overline{\nabla\Lha}_{b,k}\big),
\]
so any gradient direction shared by all members of a set cancels. What survives is the \emph{contrast}: ``move towards the better siblings and away from the worse ones, from this start and at this position''. That is the precise sense in which the optimisation is sandboxed per $(b,k)$, even after the samples are flattened into minibatches.
\end{intuition}

\subsection{The update loop}
Algorithm~1 of the paper, made explicit:
\begin{enumerate}
\item \textbf{Freeze a behaviour copy.} $\thao\leftarrow\tha$.
\item \textbf{Roll out} $B\cdot G=128$ trajectories of $N=5$ chunks with $\thao$ through Cosmos3 (UniPC) and Ctrl-World. For each of the $640$ chunks store $z$, $c=(o,g)$, the probe seeds $\Omega$, and $\Lha_{\thao}$ (computed once, under \texttt{no\_grad}).
\item \textbf{Score}: rewards \eqref{loop:eq:reward}, returns \eqref{loop:eq:return}, advantages \eqref{loop:eq:adv}. These are now constants.
\item \textbf{Optimise}: split the 640 chunks into minibatches. For each minibatch compute $\Lha_{\tha}$ with gradients, $\hat\rho$, the clipped loss $-J$ (averaged over the minibatch), backpropagate into $\tha$ only, and take \emph{one} optimiser step (AdamW with gradient-norm clipping, per the config fields). In the code the minibatch is assembled from micro-batches on each GPU with gradient accumulation and data parallelism.
\item Because $\tha$ changes after the first minibatch, $\hat\rho\neq1$ from the second minibatch on, and that is where clipping becomes active. With a single pass and few minibatches, clipping is a mild safeguard rather than the main brake.
\item Go to 1. A production run is about 50 updates; the laptop ablation in Fig.~6 shows 15.
\end{enumerate}

\begin{codenote}[(batch and ratio variants)]
\begin{itemize}[leftmargin=*]
\item The batch is fixed by a contract: global minibatch $=$ micro-batch per GPU $\times$ world size $\times$ accumulation steps, and the number of optimiser steps per update is (trajectories $\times$ chunks)$/$global minibatch. For example, a cup profile in the tests has 512 trajectories (2560 chunks) and a global minibatch of 1280 chunks, i.e.\ 2 optimiser steps per update. Some profiles thus use larger batches than the paper's $B\cdot G=128$.
\item \texttt{fpo.ratio\_granularity} $\in\{\texttt{per\_action},\texttt{per\_mc}\}$. \eqref{loop:eq:ratio} averages the loss over the $Q$ probes \emph{before} exponentiating (one ratio per chunk). \texttt{per\_mc} forms and clips one ratio per probe; a unit test checks that the two give different losses. The PPO clip may be asymmetric (\texttt{clip\_ratio\_low}, \texttt{clip\_ratio\_high}).
\item The loss logs \texttt{approx\_kl} and \texttt{clip\_fraction}, the two numbers to watch to know whether the step size is sensible.
\end{itemize}
\end{codenote}

\subsection{What is frozen, what learns}
\begin{center}
\small
\begin{tabularx}{\linewidth}{@{}l c Y@{}}
\toprule
Component & Status & Role in the update\\
\midrule
Cosmos3 video path $\thv$ & frozen & produces $\VC$ (reward target) and the video latents the action head attends to\\
Cosmos3 action head $\tha$ & \textbf{trained} & only parameters with gradients, via $\Lha$\\
Ctrl-World $\phi^0$ & frozen & produces $\VW$ and the next observation; never differentiated\\
Goal image $I^\star_g$ & fixed & from the demonstration\\
Success classifier & frozen & diagnostics only (model-space success)\\
\bottomrule
\end{tabularx}
\end{center}
A subtle point: $\VC$ is generated jointly with $a$ from the same latent. Changing $\tha$ changes the action latents during sampling, which the video latents attend to, so the plans $\VC$ can change after training even though $\thv$ is frozen. ``The policy cannot modify the predictive models'' is true for parameters, not for the plans it is scored against.

%=====================================================================
\section{One update with numbers}
%=====================================================================
\begin{center}
\small
\begin{tabularx}{\linewidth}{@{}l Y@{}}
\toprule
Step & Paper setting\\
\midrule
Groups, members, chunks & $B=8$, $G=16$, $N=5$ $\Rightarrow$ 128 rollouts, 640 chunk samples\\
Actions per chunk & $32\times10$ at 15\,Hz (about 2.1\,s); one rollout covers about 10.7\,s of imagined time\\
Model calls & 640 Cosmos3 samples (UniPC) $+$ 640 Ctrl-World predictions\\
Rewards & 640 trajectory terms, 128 goal terms (last chunk only)\\
Comparison sets & $B\cdot N=40$ sets of 16; advantages have zero mean and unit mean square in each\\
FPO evaluations & $640\times Q$ forward passes for $\Lha_{\mathrm{old}}$, the same again with gradients\\
Optimiser steps & $640/(\text{minibatch size})$, single pass\\
Discount & $\gamma=0.95$\\
Hardware & 32 GH200 (from the authors' slides); per task 928--1396 GPU-hours, 29--43 h wall clock\\
\bottomrule
\end{tabularx}
\end{center}

%=====================================================================
\section{Experimental setup (Section V.A and around)}
%=====================================================================

\subsection{Tasks, data, trials}
\begin{center}
\small
\begin{tabular}{@{}lrrrrl@{}}
\toprule
Task & Demos & Real trials$^\dagger$ & SFT & RL & Main difficulty\\
\midrule
Laptop closing & 40 & 50 & 22.0\% & 68.0\% & contact: rotate lid, keep base still\\
Duck placing & 200 (50/colour) & 60 & 95.0\% & 100.0\% & identity from instruction (4 colours)\\
Push-T & 150 & 90 & 23.3\% & 48.9\% & pose inference, pushing\\
Cup nesting & 100 & 30 & 33.3\% & 83.3\% & longest horizon, 4 grasp-and-place\\
\bottomrule
\end{tabular}

\smallskip
{\footnotesize $^\dagger$Not stated in the paper; inferred from Table~II, where $n$ is the number of successes: $11/0.22=34/0.68=50$, $57/0.95=60/1.0=60$, $21/0.233=44/0.489=90$, $10/0.333=25/0.833=30$.}
\end{center}

\begin{itemize}
\item \textbf{Robot:} UR5 (the code's kinematics are those of the original UR5), three cameras (main, wrist, side).
\item \textbf{Split and checkpoint:} demonstrations split 4:1 train/test per task; for every method the checkpoint with lowest validation loss \emph{during adaptation} is selected. For the RL stage the paper does not say which update is deployed.
\item \textbf{Where RL data comes from:} only the first frame of each demonstration (as reset state) and its final frame (as goal image). No real transitions, no success labels.
\item \textbf{Baselines:} SFT initialisation, Fast-WAM, VLA-JEPA, $\pi_{0.5}$, each with its official training recipe and execution horizon on the same demonstrations. The weak Fast-WAM/VLA-JEPA numbers are read by the authors as ``low-data regime is hard'', not as a verdict on those methods.
\item \textbf{Metric:} real-robot success rate per task, averaged with equal task weight (43.4\%$\to$75.1\%); plus completion time and control steps over successful trials (Table~II, so the two rows cover different trial sets).
\end{itemize}

\subsection{The model-space experiments}
These are not real-robot results; all their numbers come from Ctrl-World predictions scored by a frozen instruction-conditioned success classifier on the terminal frame.
\begin{itemize}
\item \textbf{Fig.~5 (gap during SFT):} at 0.5k, 5k, 10k, 15k SFT steps, 100 paired rollouts from 20 initialisations. Imagined video ``succeeds'' 18\%$\to$50\%, Ctrl-World's rollout of the same actions 6--13\%. This 41-point gap is the motivation for the trajectory reward.
\item \textbf{Fig.~6 (reward ablation, laptop):} trajectory-only, goal-only, combined; same initialisation, sampler, FPO and budget; 15 updates. Each reward rises under its own objective; model-space success peaks at updates 9--11 and then declines for all three. The authors suggest a mismatch between the optimised policy and Ctrl-World's training distribution. There is no stopping rule in the method.
\item \textbf{Figs.~7--8 (reward vs.\ predicted success):} combined reward separates predicted successes from failures clearly on laptop and cup, weakly on duck and Push-T. Both axes come from Ctrl-World, so this checks internal consistency, not physical correctness.
\end{itemize}

\subsection{What is not in the paper (and where to find it)}
\begin{center}
\small
\begin{tabularx}{\linewidth}{@{}l Y@{}}
\toprule
Missing in the paper & Where to look in the repository\\
\midrule
$\lambda_{\mathrm{traj}},\lambda_{\mathrm{goal}}$, per-view weights & reward section of the task YAML (\texttt{reward.video\_similarity.view\_weights})\\
$Q$, time distribution, shift, score form & \texttt{fpo.num\_mc\_samples}, \texttt{fpo.time\_distribution}, \texttt{fpo.training\_shift}, \texttt{fpo.score\_parameterization}\\
$\epsilon_{\mathrm{clip}}$ & \texttt{algorithm.clip\_ratio\_low/high}\\
Learning rate, Adam betas, weight decay, grad clip & \texttt{optimizer.*}\\
Minibatch size, optimiser steps per update & \texttt{batch.global\_minibatch\_chunks}, \texttt{batch.expected\_optimizer\_steps}\\
Number of updates, deployed checkpoint & \texttt{runtime.max\_updates} (50 in the duck profiles); deployment choice not documented\\
UniPC steps, guidance & \texttt{cosmos.sampler\_steps}, \texttt{cosmos.guidance}\\
Ctrl-World denoising steps, history & \texttt{ctrl\_world.runtime.*}\\
\bottomrule
\end{tabularx}
\end{center}
The \texttt{configs/} folder is not among the synced sources of this project, so the actual values have to be read on the lab machine.

%=====================================================================
\section{Where an intervention can enter the loop}
%=====================================================================
For your thesis (a referee that distrusts some Ctrl-World feedback), the derivation above shows four places to act, with different consequences:
\begin{enumerate}
\item \textbf{Reward level}, before \eqref{loop:eq:return}: replace or down-weight $r_{b,i,k}$. By \cref{loop:prop:shift}, replacing a reward by the group's mean at that chunk is neutral for that member's centred return, but \cref{loop:prop:budget} renormalises the set, so the other members are pushed harder (decision D1).
\item \textbf{Advantage level}, after \eqref{loop:eq:adv}: scale $\Ah_{b,i,k}$ by a trust weight $w\in[0,1]$. This changes the push size directly and does not touch other members; it breaks the zero-sum property of the set.
\item \textbf{Mask level}, inside \eqref{loop:eq:fpo}: drop a sample from the average (the code's \texttt{loss\_mask}). Equivalent to $w=0$ at advantage level.
\item \textbf{Set level}: skip whole sets $\mathcal S_{b,k}$ whose spread $s_{b,k}$ is not larger than the noise level (e.g.\ estimated from repeated Ctrl-World calls with different seeds), which addresses the amplification of meaningless differences.
\end{enumerate}

%=====================================================================
\section{A worked example}
\label{loop:app:toy}
%=====================================================================
One group, $G=4$ members, $N=3$ chunks, $\gamma=0.95$, $\lambda_{\mathrm{traj}}=\lambda_{\mathrm{goal}}=1$. Trajectory rewards per chunk and goal rewards:
\begin{center}
\small
\begin{tabular}{@{}c rrr r@{}}
\toprule
$i$ & $r^{\mathrm{traj}}_{i,0}$ & $r^{\mathrm{traj}}_{i,1}$ & $r^{\mathrm{traj}}_{i,2}$ & $r^{\mathrm{goal}}_i$\\
\midrule
1 & $-4.0$ & $-4.2$ & $-4.5$ & $-6.0$\\
2 & $-4.4$ & $-4.6$ & $-5.0$ & $-7.5$\\
3 & $-3.9$ & $-4.0$ & $-4.1$ & $-5.5$\\
4 & $-4.3$ & $-4.8$ & $-5.2$ & $-7.0$\\
\bottomrule
\end{tabular}
\end{center}
Combined rewards: $r_{i,2}=r^{\mathrm{traj}}_{i,2}+r^{\mathrm{goal}}_i$, e.g.\ $r_{1,2}=-10.5$. Returns (backwards) and advantages (population std, as in the paper):
\begin{center}
\small
\begin{tabular}{@{}c rrr rrr@{}}
\toprule
& \multicolumn{3}{c}{$R_{i,k}$} & \multicolumn{3}{c}{$\Ah_{i,k}$}\\
\cmidrule(lr){2-4}\cmidrule(lr){5-7}
$i$ & $k=0$ & $k=1$ & $k=2$ & $k=0$ & $k=1$ & $k=2$\\
\midrule
1 & $-17.47$ & $-14.18$ & $-10.50$ & $+0.62$ & $+0.60$ & $+0.58$\\
2 & $-20.05$ & $-16.48$ & $-12.50$ & $-1.03$ & $-1.00$ & $-1.09$\\
3 & $-16.36$ & $-13.12$ & $-9.60$ & $+1.32$ & $+1.33$ & $+1.34$\\
4 & $-19.87$ & $-16.39$ & $-12.20$ & $-0.91$ & $-0.94$ & $-0.84$\\
\midrule
$\mu_k$ & $-18.44$ & $-15.04$ & $-11.20$ & \multicolumn{3}{c}{$\sum_i\Ah=0$, $\sum_i\Ah^2=4$ per column}\\
$s_k$ & $1.57$ & $1.44$ & $1.20$ & & &\\
\bottomrule
\end{tabular}
\end{center}
Observations:
\begin{itemize}
\item The means differ by almost 7 between $k=0$ and $k=2$ only because $R_0$ sums more terms; normalising per position removes this.
\item The goal rewards spread by about $0.8$ (std), the per-chunk trajectory rewards by about $0.2$--$0.4$. The ranking ($3>1>4>2$) is essentially the goal ranking: \cref{loop:prop:mix} in action, although $\lambda_{\mathrm{traj}}=\lambda_{\mathrm{goal}}$.
\item With Bessel's correction (the code), every advantage is multiplied by $\sqrt{3/4}\approx0.87$ for $G=4$ (by $\approx0.97$ for $G=16$); the signs and the ranking do not change.
\item If all four members had rewards within $\pm0.01$ of each other (pure noise), the advantages would still have unit mean square: the same push, on noise.
\end{itemize}

